\documentclass[11pt]{letter}

\usepackage[margin=1in]{geometry}
\usepackage{amsmath}
\usepackage{amssymb}
\usepackage{hyperref}
\hypersetup{colorlinks=true,linkcolor=blue,citecolor=blue,urlcolor=blue}

\signature{Conrard Giresse Tetsassi Feugmo}
\address{Department of Chemistry and \\ Department of Physics \& Astronomy \\
University of Waterloo \\ 200 University Avenue West \\ Waterloo, ON N2L 3G1, Canada \\
\texttt{cgtetsas@uwaterloo.ca}}

\begin{document}
\setlength{\parskip}{0.4ex}

\begin{letter}{The Editors \\ \textit{Corrosion Science} \\ Elsevier}

\opening{Dear Editors,}

I am pleased to submit the manuscript \textbf{``Identifiability of a
point defect model for oxide growth on AISI~347''} for consideration as
a research paper in \textit{Corrosion Science}.

The point defect model is the standard mechanistic description of
passive film growth, and fitting it with four to six parameters to a
handful of thickness measurements is routine practice. Whether those
measurements determine the parameters is almost never tested. This
paper carries out that test for a representative case, and the answer
is that several of them are not determined, with consequences large enough to change a lifetime assessment.

The system is the duplex oxide grown on Nb-stabilized AISI~347
(X6CrNiNb18-10) in simulated boiling-water-reactor hydrothermal water
(240\,$^\circ$C, 7\,MPa, 0.4\,ppm dissolved O$_2$). I identify a reduced
five-parameter model from published depth-profile data at three exposure
times using a differentiable, physics-informed inversion, validated
throughout against an exact closed-form reference rather than against
another numerical solution.

I believe the work suits \textit{Corrosion Science} for four reasons.

\begin{itemize}
\item \textbf{It is the first identifiability characterization of a PDM
fit.} Profile likelihood, a 1000-member bootstrap, a prior-leverage test
and a residual-leverage decomposition are built into the inversion
rather than added afterward. Of the five parameters of the accepted
model, three exposure times leave one undetermined and pin a second only
to within an order of magnitude, and 95\% of the fit statistic is
carried by three chromium points, and 80\% by a single measurement. These
tools are mature in adjacent fields but, to my knowledge, have not been
brought to the PDM; the diagnosis transfers to any sparsely sampled
degradation model.

\item \textbf{The consequences are quantitative, and they point to an experiment.} The
identified kinetics fit the calibration data as well as an empirical
power law but extrapolate very differently: at ten years the power law
predicts a 3.5-to-6.3-fold thicker barrier layer than the mechanistic
535\,nm (bootstrap 95\% interval 447--704\,nm), varying
1.19-to-1.56-fold across the scanned operating envelope. A designed
exposure beyond roughly 3000\,hours would discriminate the two, a concrete experiment this analysis motivates.

\item \textbf{A first mechanistic description for this alloy, with its limits stated.} Prior characterization of AISI~347 oxide in BWR water stopped
at empirical per-element power-law fits with no shared physical
parameter set. Here the Cr-rich barrier and Fe-rich outer layers are
connected through a single kinetic description. I am explicit that the
model structure is not itself Nb-specific, that the neural inversion is
validated against an exact classical reference rather than required by
the problem, and that the constant-field closure underlying the derived
field strength does not survive a self-consistency test.

\item \textbf{A diagnosed and remedied failure mode.} I document an
inversion failure specific to sparse experimental data, the data loss satisfied while the governing equation is not, and remedy it with a
forward-solve self-consistency check, a cautionary and constructive
result for the growing use of physics-informed learning in corrosion
modeling.
\end{itemize}

The manuscript is original, has not been published elsewhere, and is not
under consideration by any other journal. I declare no competing
interests, and the work involves no human or animal subjects. All code
and the digitized data are released with provenance notes, as described
in the Data availability statement.

Thank you for considering this submission.

\closing{Sincerely,}

\end{letter}

\end{document}
